\documentclass{article} % Tạo một bản báo cáo
\usepackage[utf8]{inputenc} 
\usepackage[T5]{fontenc} % Để sử dụng Tiếng Việt 
\usepackage[fontsize=13pt]{scrextend} % Set fontsize=13pt
\usepackage[paperheight=29.7cm,paperwidth=21cm,right=2cm,left=3cm,top=2cm,bottom=2.5cm]{geometry} % Chuẩn A4, căn lề phải, trái, trên, dưới.
\usepackage{mathptmx} % Time New Roman
\usepackage{graphicx} % Thư viện chèn ảnh
\usepackage{float} % Set vị trì chèn ảnh 
\usepackage{tikz} % Thư viện tạo khung bìa 
\usetikzlibrary{calc} % Thư viện tikz
\usepackage{indentfirst} % Thư viện thụt đầu dòng 
\renewcommand{\baselinestretch}{1.2} % Giản dòng 1.2
\setlength{\parskip}{6pt} % Spacing after
\setlength{\parindent}{1cm} % Set khoảng cách thụt đầu dòng mỗi đoạn 
\usepackage{titlesec} % Thư viện để set up các kiểu chữ
\setcounter{secnumdepth}{4} % 4 Heading 
\titlespacing*{\section}{0pt}{0pt}{30pt} % Heading 1
\titleformat*{\section}{\fontsize{16pt}{0pt}\selectfont \bfseries \centering}

\titlespacing*{\subsection}{0pt}{10pt}{0pt} % Heading 2 
\titleformat*{\subsection}{\fontsize{14pt}{0pt}\selectfont \bfseries}

\titlespacing*{\subsubsection}{0pt}{10pt}{30pt} % Heading 3
\titleformat*{\subsubsection}{\fontsize{13pt}{0pt}\selectfont \bfseries \itshape}.

\titlespacing*{\paragraph}{0pt}{10pt}{0pt} % Heading 4
\titleformat*{\paragraph}{\fontsize{13pt}{0pt}\selectfont \itshape}.
\begin{document}
\begin{titlepage}
\begin{tikzpicture}[overlay,remember picture]
\draw [line width=3pt]
    ($ (current page.north west) + (3.0cm,-2.0cm) $)
    rectangle
    ($ (current page.south east) + (-2.0cm,2.5cm) $);
\draw [line width=0.5pt]
    ($ (current page.north west) + (3.1cm,-2.1cm) $)
    rectangle
    ($ (current page.south east) + (-2.1cm,2.6cm) $);
\end{tikzpicture}
\begin{center}
\vspace{-15pt} TRƯỜNG ĐẠI HỌC SÀI GÒN\\
\textbf{\fontsize{16pt}{0pt}\selectfont KHOA TOÁN ỨNG DỤNG}
\vspace{0.5cm}
\begin{figure}[H]
    \centering
    \includegraphics[width=6cm,height=6cm]{bianckh.png}
\end{figure}
\vspace{1.5cm}
\fontsize{24pt}{0pt}\selectfont BÁO CÁO\\
\vspace{12pt}
\textbf{\fontsize{32pt}{0pt}\selectfont NGHIÊN CỨU KHOA HỌC}
\vspace{1.5cm}
\end{center}
\hspace{20pt}\textbf{\fontsize{14pt}{0pt}\selectfont Đề tài:}
\begin{center}
\textbf{\fontsize{20pt}{0pt}\selectfont TỐI ƯU TUYẾN TÍNH CÓ THAM SỐ} 
\vspace{1cm}
\begin{table}[H]
    \centering
    \begin{tabular}{l l}
\fontsize{14pt}{0pt}\selectfont Sinh viên thực hiện:     & \fontsize{14pt}{0pt}\selectfont Nguyễn Thành Nam-3122480034 \vspace{6pt}\\     
                                                         & \fontsize{14pt}{0pt}\selectfont Lê Đức Anh-3122480001 \vspace{6pt}\\
\fontsize{14pt}{0pt}\selectfont Giảng viên hướng dẫn: & \fontsize{14pt}{0pt}\selectfont PGS.TS Tạ Quang Sơn
\end{tabular}
\end{table}
\vspace{1.5cm}
\fontsize{14pt}{0pt}\selectfont Thành phố Hồ Chí Minh, tháng 9 năm 2023
\end{center}
\end{titlepage}
\cleardoublepage
\thispagestyle{empty}
\begin{tikzpicture}[overlay,remember picture]
    \draw [line width=3pt]
        ($ (current page.north west) + (3.0cm,-2.0cm) $)
        rectangle
        ($ (current page.south east) + (-2.0cm,2.5cm) $);
    \draw [line width=0.5pt]
        ($ (current page.north west) + (3.1cm,-2.1cm) $)
        rectangle
        ($ (current page.south east) + (-2.1cm,2.6cm) $);
    \end{tikzpicture}
\begin{center}
    \vspace{-15pt} TRƯỜNG ĐẠI HỌC SÀI GÒN\\
    \textbf{\fontsize{16pt}{0pt}\selectfont KHOA TOÁN ỨNG DỤNG}
    \vspace{0.5cm}
\begin{figure}[H]
        \centering
        \includegraphics[width=6cm,height=6cm]{bianckh.png}
\end{figure}
    \vspace{1.5cm}
    \fontsize{24pt}{0pt}\selectfont BÁO CÁO\\
    \vspace{12pt}
    \textbf{\fontsize{32pt}{0pt}\selectfont NGHIÊN CỨU KHOA HỌC}
    \vspace{1.5cm}
\end{center}
    \hspace{20pt}\textbf{\fontsize{14pt}{0pt}\selectfont Đề tài:}
\begin{center}
    \textbf{\fontsize{20pt}{0pt}\selectfont TỐI ƯU TUYẾN TÍNH CÓ THAM SỐ} 
    \vspace{1cm}
\begin{table}[H]
        \centering
        \begin{tabular}{l l}
    \fontsize{14pt}{0pt}\selectfont Sinh viên thực hiện:     & \fontsize{14pt}{0pt}\selectfont Nguyễn Thành Nam - 3122480034 \vspace{6pt}\\     
                                                             & \fontsize{14pt}{0pt}\selectfont Lê Đức Anh - 3122480001 \vspace{6pt}\\
    \fontsize{14pt}{0pt}\selectfont Giảng viên hướng dẫn: & \fontsize{14pt}{0pt}\selectfont PGS.TS Tạ Quang Sơn
\end{tabular}
\end{table}
    \vspace{1.5cm}
    \fontsize{14pt}{0pt}\selectfont Thành phố Hồ Chí Minh, tháng 9 năm 2023
\end{center}
\cleardoublepage
\section*{LỜI NÓI ĐẦU}
\thispagestyle{empty}
Phần này trình bày một cách rất khái quát (khoảng 1-2 trang) về bối cảnh hình thành và mục đích của báo cáo. Lời cảm ơn với những tổ chức và cá nhân góp phần trong việc hoàn thiện báo cáo (nếu có) nên đặt ở cuối mục này.
\cleardoublepage
\section*{LỜI CAM ĐOAN}
\thispagestyle{empty}
Chúng em gồm: NGUYỄN THÀNH NAM, mã số sinh viên 3122480034, sinh viên lớp DTU1221, khóa 2022-2026, LÊ ĐỨC ANH, mã số sinh viên 31222480001, sinh viên lớp DTU1221, khóa 2022-2026, người hướng dẫn là PGS.TS.TẠ QUANG SƠN. Chúng em xin cam đoan toàn bộ nội dung được trình bày trong báo cáo Tối Ưu Tuyến Tính có Tham Số là kết quả quá trình tìm hiểu và nghiên cứu của chúng em. Các dữ liệu được nêu trong báo cáo là hoàn toàn chung thực, phản ánh đúng kết quả đo đạc thực tế. Mọi thông tin trích dẫn đều tuân thủ các quy định và sỡ hữu trí tuệ: các tài liệu tham khảo được liệt kê rõ ràng. Chúng em xin chịu hoàn toàn trách nhiệm với những nội dung được viết trong báo cáo này.

\vspace{6pt}
\hspace{4.5cm}Thành phố Hồ Chí Minh, ngày 14 tháng 9 năm 2023 

\hspace{7.5cm}\textbf{Người cam đoan}

\vspace{2cm}
\hspace{4cm}\textbf{LÊ ĐỨC ANH} \hspace{3.5cm}\textbf{NGUYỄN THÀNH NAM} 
\cleardoublepage

\addtocontents{toc}{\protect\thispagestyle{empty}}
\tableofcontents % Tạo mục lục tự động 
\thispagestyle{empty}
\cleardoublepage

\pagenumbering{roman} % Đánh số thứ tự la mã 
\section*{DANH MỤC KÝ HIỆU VÀ CHỮ VIẾT TẮT}
\addcontentsline{toc}{section}{\numberline {} DANH MỤC KÝ HIỆU VÀ CHỮ VIẾT TẮT}
\cleardoublepage

\listoffigures % Tạo danh mục hình vẽ tự động
\addcontentsline{toc}{section}{\numberline {} DANH MỤC HÌNH VỄ}
\cleardoublepage

\listoftables % Tạo danh mục hình vẽ tự động
\addcontentsline{toc}{section}{\numberline {} DANH MỤC BẢNG BIỂU}
\cleardoublepage
\section*{TÓM TẮT BÁO CÁO}
\addcontentsline{toc}{section}{\numberline {} TÓM TẮT BÁO CÁO}
Phần này trình bày những mục đích và các kết luận quan trọng nhất của báo cáo bằng cả Tiếng Việt và Tiếng Anh. 
\cleardoublepage
 
\pagenumbering{arabic} % Đánh số thứ tự 1,2,3,... 
\section*{CHƯƠNG 1. CHƯƠNG MỞ ĐẦU}
\addcontentsline{toc}{section}{\numberline {} CHƯƠNG 1. CHƯƠNG MỞ ĐẦU}
Phần mở đầu giới thiệu vấn đề mà báo cáo cần giải quyết, mô tả được các phương pháp hiện có để giải quyết vấn đề,
trình bày mục đích của báo cáo song song với việc giới hạn phạm vi của vấn đề mà báo cáo sẽ tập trung giải quyết.
Phần này cũng sẽ giới thiệu tóm tắt cấu trúc báo cáo và nội dung tương ứng của các phần sẽ lần lượt được trình bày
ở các chương tiếp theo.

Nội dung chính của một báo cáo thường bao gồm: 
\begin{itemize}
    \item Phần mở đầu giới thiệu đề tài.
    \item Một chương giới thiệu cơ sở lý thuyết.
    \item Một hoặc nhiều chương trình bày các vấn đề về tính toán và thiết kế.
    \item Một chương mô tả các thí nghiệm và kết quả thu được.
\end{itemize}
\newpage

\section*{CHƯƠNG 2. CƠ SỞ LÝ THUYẾT}
\addcontentsline{toc}{section}{\numberline {} CHƯƠNG 2. CƠ SỞ LÝ THUYẾT}
\setcounter{section}{2}
\setcounter{subsection}{0}
Mỗi chương sẽ bắt đầu bằng một đoạn giới thiệu các phần chính được trình bày  trong chương đó, dài khoảng từ 5 đến 10
dòng và kết thúc bằng một đoạn tóm tắt các kết luận chính của chương. Chú ý phân bố chiều dài của mỗi chương cho cân
đối và hợp lý.
\subsection{Một số lưu ý khi trình bày báo cáo}
Sau đây là một vài chú ý khi làm báo cáo cần nhớ:
\subsubsection{Nộp báo cáo}
\vspace{-35pt}
Sinh viên (hoặc nhóm sinh viên tối đa 3 thành viên làm chung một đề tài) nộp 2 quyển báo cáo nghiên cứu khoa học tại
văn phòng bộ môn của giảng viên hướng dẫn trước ngày phản biện ít nhất 1 tuần. Một quyển báo cáo cần có các đặc điểm
sau:
\begin{itemize}
    \item Được \textbf{in hai mặt} nhằm tiết kiệm không gian lưu trữ.
    \item Đóng bìa mềm, bên ngoài là bóng kính.
    \item Số trang: 50 - 150 trang, không kể phần phục lục.
    \item Phải có chữ ký của sinh viên sau LỜI CAM ĐOAN và của giảng viên hướng dẫn.
\end{itemize}
\subsubsection{Phụ lục}
\vspace{-35pt}
Phụ lục (nếu có) chứa các thông tin có liên quan đến báo cáo nhưng nếu để ở trong phần chính sẽ gây rườm rà. Thông
thường các chi tiết được để trong phần phụ lục là: kết quả thô (chưa qua xử lý), mã nguồn phần mềm, thông số kỹ thuật
chi tiết của linh kiện, hình ảnh minh họa thêm,...vv.

\subsubsection{Tài liệu tham khảo}
\vspace{-35pt}
\paragraph{Cách liệt kê}\mbox{}

Áp dụng cách liệt kê theo quy định của IEEE. Theo đó, tài liệu tham khảo được đánh số thứ tự trong ngoặc vuông. Thứ
tự liệt kê là thứ tự xuất hiện của tài liệu tham khảo được trích dẫn trong báo cáo. Tài liệu tham khảo đã liệt kê
bắt buộc phải được trích dẫn trong phần nội dung của báo cáo. Tài liệu tham khảo cần có nguồn gốc rõ ràng và phải
từ nguồn đáng tin cậy. Hạn chế trích dẫn tài liệu tham khảo từ các website, từ wikipedia.
\paragraph{Các loại tài liệu tham khảo}\mbox{}

Các nguồn tài liệu tham khảo chính là sách, bài báo trong các hội nghị khoa học và các tài liệu tham khảo khác trên
internet.

 \subsubsection{Đánh số phương trình}
 \vspace{-30pt}
 Phương trình được đánh số theo số của chương, như hình vẽ và bảng biểu.
 \subsubsection{Đánh số định nghĩa, định lý, hệ quả}
 \vspace{-30pt}
 Các định nghĩa, định lý, hệ quả sẽ được đánh số theo số của chương và được sử dụng chung một chỉ số. Ví dụ trong
 chương 3, các định nghĩa, định lý, hệ quả sẽ được đánh số theo thứ tự như sau: Định lý 3.1, Định nghĩa 3.2, Hệ quả
 3.3, Định lý 3.4,...
\end{document}